\section{Related Work}
\label{sec:relatedwork}

\subsection{Exposure, Preparatory, and Deformable-Object Manipulation}

Robotic manipulation often includes actions whose purpose is to establish conditions for a later operation. Manipulation-based object search reasons about visibility and reachability while moving occluding objects~\cite{dogar2014search}; mechanical search combines repeated perception and pushing, suction, or grasping to retrieve a target hidden in clutter~\cite{danielczuk2019mechanical}. Preparatory-manipulation planners instead regrasp, place, or hand over an object to reach a configuration suitable for later use~\cite{wan2020preparatory}. Recent BiPreManip learns anticipatory affordances for asymmetric bimanual tasks in which one arm reconfigures an object for the other arm's goal-directed operation~\cite{shen2026bipremanip}. These works establish the importance of actions that make later manipulation feasible. Mechanical search primarily removes rigid clutter to find or retrieve a target, while preparatory manipulation primarily changes the configuration of the manipulated object. Our setting combines physical exposure, target execution, and restoration under a deformable occluder, and evaluates whether a continuous visuomotor policy preserves the changing conditions connecting these operations.

Occlusion can also be addressed by changing sensing rather than changing the scene. Active-vision systems move a camera to improve task-specific viewpoints for bimanual imitation learning~\cite{chuang2025activevision}. In contrast, exposure in our tasks is itself a contact-rich manipulation objective: moving the occluder changes both what can be observed and what motions are physically feasible. Surgical tissue retraction provides a particularly relevant example. Prior systems study autonomous retraction with engineered perception and control~\cite{attanasio2020retraction}, learning from demonstrations~\cite{pore2021retraction}, or adaptive visual control for exposing a region of interest~\cite{liu2026exposure}. These approaches focus on attaining and maintaining exposure. We study the broader dependency between exposure and the subsequent target operation, including restoration as part of the complete-task criterion.

Deformable-object manipulation is challenging because cloth, rope, tissue, and bags have high-dimensional configurations, complex contact dynamics, and perceptual ambiguity~\cite{jimenez2012deformable,sanchez2018deformable}. Model-based methods encode object geometry and mechanics, whereas data-driven methods learn state estimates, dynamics, or policies from interaction data. Representative learning-based systems acquire visual policies for cloth and rope manipulation without demonstrations~\cite{wu2020deformable}, combine goal-conditioned local and global actions for deformable objects~\cite{kim2024dualaction}, or learn structure-of-interest dynamics for bimanual bag manipulation~\cite{zhou2025bag}. Rather than reconstructing or predicting the complete deformable state, our method selects task descriptors that summarize visibility and spatial relationships needed for action selection. This is a deliberate abstraction: RGB and proprioception remain available for contact and deformation cues not represented by those descriptors.

\subsection{Visuomotor Imitation for Sequential and Bimanual Manipulation}

Visuomotor imitation learning maps sensory observations and robot state to actions without requiring an explicit contact model. Modern policies often predict temporally extended commands. ACT uses a transformer to generate action chunks for fine-grained bimanual manipulation~\cite{zhao2023act}, whereas Diffusion Policy models multimodal action sequences through conditional denoising~\cite{chi2023diffusion}. These policy families support coordinated and multistage behaviors in Mobile ALOHA and UMI~\cite{fu2024mobilealoha,chi2024umi}. Recent work also improves bimanual data efficiency by extracting reusable motion patterns from human videos and learning a bimanual diffusion policy~\cite{zhou2025yoto}. Such methods provide strong action-generation backbones, but action prediction alone does not specify which consequences of earlier interactions should remain explicit in the policy representation.

Long-horizon manipulation can instead be organized with explicit task structure. Neural task programs and learned sensorimotor primitives decompose demonstrations and select skills using visual state, object relations, or action history~\cite{xu2018ntp,liang2022sensorimotor}. C-LEARN represents multistep demonstrations as ordered keyframes with geometric constraints~\cite{perez2017clearn}, while Relay Policy Learning uses goal-conditioned hierarchical policies initialized from unstructured demonstrations and refined with reinforcement learning~\cite{gupta2020relay}. For physical multistage tasks, hierarchical imitation has been used to select and retry primitives for cable routing~\cite{luo2024cable}, and dual-action imitation combines reactive local commands with generated global trajectories for deformable-object manipulation~\cite{kim2024dualaction}. These approaches address temporal abstraction, skill sequencing, constraint satisfaction, or trajectory generation. Our work keeps a single action-chunk policy and fixed offline demonstrations, and instead studies whether the representation used by that policy retains achieved exposure and interaction geometry. It consequently requires neither primitive annotations nor an explicit phase, subgoal, or high-level skill selector.

\subsection{Task-Relevant Representations for Imitation Learning}

Behavior cloning can fit demonstrated actions using correlations that are predictive in the dataset but unreliable during closed-loop execution. Causal-confusion analyses show that an imitation learner can select the wrong causal explanation for expert behavior~\cite{dehaan2019}. With observation histories, policies can similarly copy preceding actions rather than use environmental evidence; adversarial history dropout and residual-action prediction were developed to counter this shortcut~\cite{wen2020copycat,chuang2022rap}. These studies concern temporal or causal shortcuts rather than exposure geometry specifically, but they motivate our central distinction: small action-prediction error does not guarantee that the learned representation preserves the physical information on which an appropriate action depends.

Structured visual representations reduce the burden of extracting such information from pixels. Transporter Networks encode pick-and-place geometry through feature transport~\cite{zeng2021transporter}; VIOLA uses object proposals and attention to emphasize regions useful for action generation~\cite{zhu2023viola}; and GROOT learns object-centric 3-D representations for generalization across scenes and instances~\cite{zhu2023groot}. SAM-E adapts a visual foundation model to extract task-relevant features for sequence imitation~\cite{zhang2024same}. More recently, task-aware semantic--geometric observation interfaces use open-vocabulary segmentation and depth to canonicalize policy inputs under appearance shifts~\cite{ding2026choose}. These methods demonstrate the value of spatial and object-centric priors. Our method further constrains internal action-generating queries so that selected task descriptors are recoverable; the ablations therefore separate the effect of supplying semantic observations from the effect of supervising what the joint policy representation retains.

Task awareness has also been introduced through learning objectives and additional teaching signals. Belief-module imitation learning jointly trains a history representation and an adversarial policy with a task-aware imitation loss~\cite{gangwani2020belief}. TRAIL restricts an adversarial-imitation discriminator to avoid spurious task-irrelevant visual features~\cite{zolna2021trail}, while information-bottleneck behavior cloning suppresses latent redundancy while preserving features useful for manipulation~\cite{bai2025ibbc}. CIVIL asks demonstrators to identify decision-relevant objects and positions through language and physical markers, then learns a compact causal visual representation with the policy~\cite{dai2026civil}. TARAD uses language-model planning and visual affordances to generate demonstrations for a multitask diffusion policy~\cite{hu2025tarad}. Thus, \emph{task-aware imitation learning} is not introduced here as a new general term. In our setting it has a narrower operational definition: mask-derived descriptors of target visibility, tool--target displacement, and target--goal distance are supervised directly from policy query tokens and trained jointly with action prediction. They are neither task identifiers nor reward or discriminator features; they state what physical information should remain readable from the representation producing coordinated bimanual commands.

Our descriptor targets depend on scalable semantic supervision. Promptable video segmentation models such as SAM2 propagate sparse prompts through video~\cite{ravi2024sam2}, and semi-supervised segmentation methods expand sparse labels using consistency-based pseudo-labeling~\cite{yang2023unimatch}. Because persistent segmentation errors can yield plausible yet incorrect geometric descriptors, our workflow combines reviewed bidirectional masks with video-level temporal and structural screening before cumulative self-training. This annotation procedure supplies supervision to the policy; it is not itself assumed to improve control, and its effect is evaluated through the downstream factorial ablation.

\subsection{Distribution Shift and Task-Level Analysis}

Behavior cloning can accumulate errors when a learned policy visits states that are poorly represented in demonstrations. DAgger addresses this distribution shift by querying the expert on learner-visited states~\cite{ross2011}, whereas DART perturbs demonstrations to collect behavior around likely execution errors~\cite{laskey2017dart}. These approaches improve deployment coverage but require additional expert interaction or a modified collection process. Our method retains a fixed demonstration set and examines how task-directed representation supervision changes the policy learned from it; it does not claim to eliminate unsupported deployment histories.

Several analyses connect local imitation decisions to longer-term performance. AggreVaTe and value aggregation evaluate actions through downstream task cost~\cite{ross2014aggrevate,cheng2018value}, while stability-constrained imitation learning and TaSIL control error propagation under regularity assumptions on the system and expert~\cite{tu2022,pfrommer2022tasil}. Representation and domain-adaptation analyses further show that distribution alignment can discard information needed for prediction unless support and recoverability conditions are preserved~\cite{johansson2019support,zhao2019invariant}. Our analysis complements these results. Under explicit coverage and action-sensitivity conditions, it connects representation-level imitation error to failure of the complete exposure-and-execution task and gives expert-action reconstruction from task descriptors as a sufficient route for controlling the demonstration-to-deployment error discrepancy. These are conditional design guarantees, not an assertion that semantic supervision alone ensures task success.
